\DTMsavedate{mydate}{2022-06-30}
\maketitle{Analogue Communications Engineering}{Electronic Engineering}{\DTMusedate{mydate}}

% ----------------------------------------------------- %
% COURSE INFO
\subsection*{Course Information}\label{course:tee3121}
\vspace{0ex}%
\noindent\parbox{0.5\textwidth}{%
    \noindent\begin{tabular}{@{}ll}
                 \textsf{Course:}          & 	Analogue Communications Engineering  \\
                 \textsf{Course Code:}         & 	TEE 3121    \\
                 \textsf{Credits:}    & 	10
    \end{tabular}}
\vspace{2ex}\hrule\vspace{2ex}




% ----------------------------------------------------- %
% INSTRUCTOR INFO
\subsection*{Lecturer Information}
\noindent\parbox{0.5\textwidth}{%
    \noindent\begin{tabular}{@{}ll}
                 \textsf{Name:}         &     Magripa Nleya         \\
%\textsf{Phone:}        &    \phone          \\
\textsf{Email:}        &    magripa.nleya@nust.ac.zw \\
\textsf{Qualifications:}        &  MSc in Telecommunication Engineering
    \end{tabular}}
\vspace{2ex}\hrule\vspace{2ex}

% ----------------------------------------------------- %
% COURSE SYNOPSIS
\subsection*{Course Synopsis}
This is an introduction to analogue telecommunications covering amplitude modulation, single-side band modulation, angle modulation, frequency division multiplexing, propagation effects, demodulation.
It also covers determination of the signal-to-noise ratio in AM and FM systems.
There is design of small-signal HF amplifiers, mixers, oscillators and detectors, mixer theory and spectral analysis, noise generation in electronic circuits and devices.
\vspace{2ex}\hrule\vspace{2ex}

% ----------------------------------------------------- %
% COURSE DESCRIPTION
\subsection*{Learning Objectives}
The students will be able to:
\begin{outline}
    \1 understand telecommunication principles.
    \1 learn about mobile communication systems.
    \1 to have an understanding of modulation and multiplexing methods in telecommunication systems.
\end{outline}
\vspace{2ex} \hrule\vspace{2ex}

% ----------------------------------------------------- %
% COURSE TOPICS
\subsection*{Topics}
\begin{outline}
    \1 Communication systems
    \1 Analogue cellular systems
    \1 Amplitude Modulation
    \1 Angle Modulation
    \1 Transmission and propagation
    \1 Multiplexing
    \1 Noise in communication systems
\end{outline}
\vspace{2ex} \hrule\vspace{2ex}

% ----------------------------------------------------- %
% Students Assenssment
\subsection*{Assessment of Students}
\begin{outline}
    \1 Students are assessed through assignments, laboratory and written in-class tests which shall form the continuous assessment mark.
    \1 There shall be final examination at the end of the semester.
\end{outline}
\vspace{2ex}\hrule\vspace{2ex}

% ----------------------------------------------------- %
% Grade Evaluation
\subsection*{Course Evaluation and Grading Policies}
Grades are based on:
Continuous assessment (\qty{25}{\percent}),
Final Exam (\qty{75}{\percent})
\vspace{2ex}\hrule\vspace{2ex}

% ----------------------------------------------------- %
% EQUIPMENT
% https://sites.udel.edu/computing-purchases/personal-specs/
% \subsection*{Essential Equipment}

% \vspace{2ex}\hrule\vspace{2ex}

% Policies and Other information
\subsection*{Policies and Other Information}
% Conduct
\subsection*{Key Text}
\begin{itemize}
    \item Communications engineering principles by Ifiok Otung
    \item Lecture notes
\end{itemize}
\vspace{2ex}\hrule
\vspace{2ex}\hrule\vspace{2ex}